\documentclass[10pt,a4paper]{amsart}
\usepackage[margin=19mm]{geometry}
\usepackage{amsmath,amssymb,mathtools}
\usepackage[scr=rsfs,cal=euler]{mathalfa}
\usepackage{xfrac}
\usepackage[unicode,hidelinks,bookmarks=false]{hyperref}
\newcommand{\ie}{i.e.\ }
\newcommand{\cf}{cf.\ }
\newcommand{\resp}{respectively\ }
\newcommand{\Subsection}{\subsection}
\providecommand{\nobreakdash}{\nobreak\hbox{-}}
\setlength{\emergencystretch}{2em}
\multlinegap0pt
\usepackage{bm}
\usepackage{bbm}
\usepackage{dsfont}
\usepackage{xspace}
\usepackage{cancel}
\usepackage{mathtools}
\usepackage{amsthm}
\makeatletter
\@ifundefined{lemma}{\newtheorem{lemma}{Lemma}}{}
\@ifundefined{proposition}{\newtheorem{proposition}{Proposition}}{}
\makeatother
\def\RR{{\mathcal{R}}}
\def\R{{\mathbb R}}
\def\a{\alpha}
\def\d{\delta}
\def\e{\varepsilon}
\def\g{\gamma}
\def\ti{\tilde}
\def\to{\rightarrow}
\def\uv{\underline{v}_{\d}}
\title[A maximal estimate for the non-linear Fourier transform]{A maximal estimate for the non-linear Fourier transform}
\author{Alexei Poltoratski}
\date{Source: arXiv:2608.24606v1 (CC BY 4.0)}
\begin{document}
\maketitle

\noindent\emph{Source and licence.} A maximal estimate for the non-linear Fourier transform. Version arXiv:2608.24606v1 (CC BY 4.0), distributed under the Creative Commons Attribution 4.0 International licence, as recorded in the arXiv licence metadata for this exact version and shown on the abstract page https://arxiv.org/abs/2608.24606v1. Full rights notice: licenses/arxiv-2608.24606-CC-BY-4.0.txt.\par\smallskip

\noindent\textit{Editorial scope.} Counted unit: the Lemma whose source label is \texttt{lemSweep} in the article (source0.tex lines 638--653, source label \texttt{lemSweep}), delivered together with the article's own complete original proof of it (source0.tex lines 655--715, one \texttt{proof} environment of 61 lines). No mathematics is written, repaired or re-derived: statement and proof are the article's own text line for line. Required context retained uncounted: propU3 (lines 387--396); eqEpsU3 (lines 442--445); lemRep (lines 487--500); lemChi (lines 517--521); lemPair (lines 530--552); eqDepthMatch (lines 540--543); eqDeltaPin (lines 546--549); fin11 (lines 620--622). Every definition, display and label of those lines is retained unchanged. Omitted from this package and not redistributed: 1--386, 397--441, 446--486, 501--516, 522--529, 544--545, 550--619, 623--637, 654--654, 716--1260. Nothing inside a retained excerpt is omitted or altered: every retained body is delivered exactly as the article's lines are written, so no omission receipt of any kind accompanies this package. Citations: the retained excerpts carry no citation command at all, so no bibliography entry of the article is transcribed and no reference is redistributed. Reference adaptation: none was needed and none was made; every reference inside the delivered unit resolves inside it, so no unresolved reference or citation is produced. Printed slips kept literal: no printed word, symbol, hypothesis or number of the counted statement or of its proof was altered. Layout and class: the article's own class is \texttt{amsart} with packages including amssymb, xcolor, hyperref; the wrapper is the lane's pinned \texttt{amsart} editorial wrapper at 10pt on A4, which restates the theorem-like environments the retained text needs and adds the article's own macro definitions that the retained text needs (\texttt{\textbackslash R}, \texttt{\textbackslash RR}, \texttt{\textbackslash a}, \texttt{\textbackslash d}, \texttt{\textbackslash e}, \texttt{\textbackslash g}, \texttt{\textbackslash ti}, \texttt{\textbackslash to}, \texttt{\textbackslash uv}), with extra wrapper packages (bm, bbm, dsfont, xspace, cancel, mathtools). The delivered numbering is the unit's own. Assets: no retained span contains a figure environment, a graphics-inclusion command, a PDF-page inclusion, a table or a TikZ picture; no image file is copied into this package, the package declares no asset dependency and no image file is redistributed with it. Licence: version arXiv:2608.24606v1 of this article is distributed under the Creative Commons Attribution 4.0 International licence, as recorded in the arXiv licence metadata for this exact version and shown on the abstract page https://arxiv.org/abs/2608.24606v1. A full rights notice is delivered at licenses/arxiv-2608.24606-CC-BY-4.0.txt. Characters: every retained excerpt is pure ASCII, so the delivered engine, XeLaTeX with its default Latin Modern text font, reports no missing character.
\begin{proposition}\label{propU3}
There are absolute constants $K_9$, $K_{11}$ such that the following holds on $\RR(T,\eta_0,C,\d)$ for $t\geq T$ and $c\in\{C,2C,4C,8C\}$, provided that $C\geq4\pi$ and $4\d\leq\eta_0$. (Both conditions hold in all applications below: there $C\geq C_1(\eta_0)\geq4\pi$, with $C_1$ of Lemma \ref{lemPair}, and the thresholds imposed on $\d$ are always taken smaller than $\eta_0/4$.)

1) {\rm(lattice rigidity)} If $t\notin T_0(s,c)$, then every zero of $E(t,\cdot)$ in $Q(s,8C^\sharp/t)$ has depth at least $c-1$.

2) {\rm(resonance-free models)} If $t\notin T_0(s,c)$ and $v_{\min}\in[c-1,\infty]$ denotes the smallest depth of the zeros of $E(t,\cdot)$ in $Q(s,8C^\sharp/t)$, then there is a unimodular constant $\a$ such that
$$\left|E(t,x)-\frac{\a}{\sqrt w}\,e^{-itx}\right|\ \leq\ K_{11}\left(\d+e^{-2v_{\min}}\right)\eta_0^{-1}
\qquad\text{ for all }x\in\R\cap Q(s,c/t).$$
In particular, if no zeros of $E(t,\cdot)$ in $Q(s,8C^\sharp/t)$ have depth $\leq(C^\sharp-C)/2$, then (U3) holds at $(s,t)$ on the box $Q(s,c/t)$ with $\e=K_{11}\,\d\,\eta_0^{-1}$. The same statements hold for $\ti E$.
\end{proposition}

\begin{equation}
\epsilon_c:=K_{11}\left(\d+e^{-2(c-1)}\right)\eta_0^{-1}
\label{eqEpsU3}
\end{equation}

\begin{lemma}\label{lemRep}
Let $\zeta$ be a zero of $E(t,\cdot)$ in $Q(s,8C/t)$, let $v$ be its depth, and suppose $v\geq\uv$. Then there is a unimodular constant $\a=\a(s,t)$ such that
\begin{equation}
\sup_{z\in Q(s,8C/t)}\left|E(t,z)-\a\,\frac{\g(v)}{\sqrt{w}}\,\sin[t(z-\zeta)]\right|\ \leq\ \vartheta\,\g(v),
\quad
\vartheta:=K_5\,\d^{1/2}e^{16C}\eta_0^{-3/2},
\label{eqRepE}
\end{equation}
with an absolute constant $K_5$. The same holds for the Dirichlet family: if $\ti\zeta_0$ is a zero of $\ti E(t,\cdot)$ in $Q(s,8C/t)$ of depth $\ti v\geq\uv$, then, with $\ti\zeta:=\ti\zeta_0-\pi/(2t)$ and a unimodular $\ti\a$,
\begin{equation}
\sup_{z\in Q(s,8C/t)}\left|\ti E(t,z)-\ti\a\,\frac{\g(\ti v)}{\sqrt{\ti w}}\,\cos[t(z-\ti\zeta)]\right|\ \leq\ \vartheta\,\g(\ti v).
\label{eqRepEt}
\end{equation}
\end{lemma}

\begin{lemma}\label{lemChi}
Suppose that $E(t,\cdot)$ has no zeros in $Q(s,c/t)$ and $\ti E(t,\cdot)$ has no zeros in $Q(s,c'/t)$, for some $c,c'\in\{C,2C,4C,8C\}$, and let $\a,\ti\a$ be the unimodular constants of the exponential models of Proposition \ref{propU3}, 2). Then $\chi:=\a\bar{\ti\a}$ satisfies
$$\left|\Im\chi-\sqrt{w\ti w}\right|\ \leq\ K\,\epsilon_c\!\left(\min(c,c')\right)\eta_0^{-1/2},$$
with $\epsilon_c$ defined in \eqref{eqEpsU3} and an absolute constant $K$.
\end{lemma}

\begin{lemma}[pairing]\label{lemPair}
There exist absolute constants $C_0\geq4\pi$, $K_6$ and $K_{12}$ with the following properties.

The constants $C_0$ and $K_{12}$ determine the least admissible scale: given $\eta_0\in(0,\tfrac12]$, we put
$$C_1(\eta_0)\ :=\ C_0+K_{12}\left(1+\log(1/\eta_0)\right)$$
and assume $C\geq C_1(\eta_0)$; the constant $K_6$ appears in the estimates \eqref{eqDepthMatch} and \eqref{eqDeltaPin} below. For every such $C$ and $\eta_0$ there is a threshold $\d_0(C,\eta_0)>0$ such that the following holds whenever $\d\leq\d_0(C,\eta_0)$; here and in the similar statements below the thresholds imposed on $\d$ are always taken smaller than $\eta_0/4$, as required by Proposition \ref{propU3}.

Let $s\in\RR(T,\eta_0,C,\d)$ and $t\geq T$, and suppose that the pair $(s,t)$ is not $\d$-shallow. Let $c$ be one of the values $C$, $2C$, $4C$, and suppose that $E(t,\cdot)$ has a zero $\zeta=x_0-iy_0$ in $Q(s,c/t)$; let $v=ty_0$ be its depth (since $(s,t)$ is not $\d$-shallow, necessarily $v\geq\uv$). Then:

1) $\ti E(t,\cdot)$ has zeros in $Q(s,2c/t)$; let $\ti\zeta_0$ be one nearest to $x_0$ and $\ti v$ its depth. Then
\begin{equation}
\left|\frac{\sinh[2\ti v]}{\sinh[2v]}-1\right|\ \leq\ K_6\,\vartheta\,\eta_0^{-3},
\label{eqDepthMatch}
\end{equation}

2) with $\d_p:=t\,(\Re\ti\zeta-x_0)$, where $\ti\zeta=\ti\zeta_0-\pi/(2t)$,
\begin{equation}
\left|\cos^2\d_p-w\ti w\right|\ \leq\ K_6\,\vartheta\,\eta_0^{-3}.
\label{eqDeltaPin}
\end{equation}

The same holds with the roles of $E$ and $\ti E$ exchanged.
\end{lemma}

\begin{equation}
b_{t_1 \to t_2}(z)=\frac {e^{i(t_2-t_1)z}}{2i}\left( \ti E(t_2,z)E(t_1,z)-E(t_2,z)\ti E(t_1,z)\right) .
\label{fin11}\end{equation}

\begin{lemma}[two-endpoint estimate]\label{lemSweep}
After enlarging the absolute constants $C_0$ and $K_{12}$ of Lemma \ref{lemPair} and decreasing the thresholds $\d_0(C,\eta_0)$, if necessary, the following holds under the standing assumptions and the endpoint configurations described above:
\begin{equation}
\left\langle |b_{t\to2t}|^2\right\rangle_{J(s)}\ \geq\ \frac{\cosh[2\nu]}{2^8\,S},
\qquad
\sup_{J(s)}|b_{t\to2t}|^2\ \leq\ \frac{2^5\cosh[2\nu]\,e^{8C}}{\eta\,S},
\label{eqSweepMoments}
\end{equation}
where $\langle\cdot\rangle_{J(s)}$ denotes the mean over $J(s)$, and in the free case $v_2=\infty$ the quantity $\cosh[2\nu]/S$ is understood as its limit $e^{-2v_1}/(2\sinh[2v_1])$.

Consequently, on a subset of $J(s)$ of measure at least $2^{-14}\,\eta_0^2\,e^{-8C}\,|J(s)|$,
\begin{equation}
\log|a_{t\to2t}|\ \geq\ \Phi(v_1,C):=\frac12\log\left(1+\frac{2^{-13}e^{-8C}}{v_1}\right).
\label{eqSweepLevel}
\end{equation}
\end{lemma}

\begin{proof}
Let $p\in J(s)$, so that in \eqref{fin11} $E^\#=\bar E$ and $\ti E^\#=\bar{\ti E}$. Write
$$\theta_k=t_k(p-\Re\zeta_k),\quad S_k=\sin(\theta_k+iv_k),\quad C_k=\cos(\theta_k-\d_k+iv_k),\quad k=1,2.$$
By Lemmas \ref{lemRep} and \ref{lemPair}, at each endpoint with presence,
\begin{equation}
\begin{gathered}
E(t_k,p)=\a_k\,\frac{\g(v_k)}{\sqrt{w}}\,S_k+O(\vartheta')\g(v_k),\\
\ti E(t_k,p)=\ti\a_k\,\frac{\g(v_k)}{\sqrt{\ti w}}\,C_k+O(\vartheta')\g(v_k),
\end{gathered}
\label{eqEndpApp}
\end{equation}
where $\a_k,\ti\a_k$ are unimodular constants and
$$\vartheta'\ :=\ K_7\left(\vartheta\eta_0^{-3}\right)^{1/2}e^{4C}:$$
here the Dirichlet depth has been replaced by $v_k$ using \eqref{eqDepthMatch}, which perturbs the depth by $|\nu_k|\leq K(\vartheta\eta_0^{-3})^{1/2}$ and hence the factor $\g(\ti v_k)\cos[\cdot+i\ti v_k]$ by at most $K(\vartheta\eta_0^{-3})^{1/2}e^{4C}\g(v_k)$ (the derivatives of the model in the depth are bounded by $e^{4C}$ times its amplitude on the box), and the Dirichlet center has been written as $\Re\zeta_k+\d_k/t_k$. If $t_2$ is in the free case, \eqref{eqEndpApp} holds at $t_2$ with $S_2$, $C_2$, $\g(v_2)$ replaced by $ie^{-i\theta_2}$, $e^{-i(\theta_2-\d_2)}$, $1$, where $\theta_2=t_2(p-x_2)$ with an arbitrary reference point $x_2$ (absorbed into the unimodular constants), by Proposition \ref{propU3}, 2) and the definition of $\d_2$; the error terms of \eqref{eqEndpApp} at the far endpoint are then $O\left(\vartheta'+\epsilon_{2C}\right)$ in place of $O(\vartheta')\g(v_2)$, and \eqref{eqRbound} below holds with $\vartheta'$ replaced by $\vartheta'+K\epsilon_{2C}$ and $\g(v_2)$ by $1$. Below the two cases are treated together, the free case corresponding to $v_2=\infty$.

Multiplying the relations \eqref{eqEndpApp} and inserting them into \eqref{fin11}, we obtain
\begin{equation}
2ie^{-i(t_2-t_1)p}\,b_{t\to2t}(p)=\frac{\g(v_1)\g(v_2)}{\sqrt{w\ti w}}\,
\left[\chi_2' \,C_2S_1-\chi_1'\, S_2C_1\right]+R(p),
\label{eqModelB}
\end{equation}
where $\chi_2'=\a_1\ti\a_2$ and $\chi_1'=\a_2\ti\a_1$ are unimodular constants, on which no assumption is made, and $R(p)$, the sum of the products containing at least one error term of \eqref{eqEndpApp}, satisfies, since $|S_k|,|C_k|\leq\cosh v_k\leq e^{4C}$ on $J(s)$,
\begin{equation}
|R(p)|\ \leq\ K\,\vartheta'\,e^{4C}\eta_0^{-1/2}\,\g(v_1)\g(v_2),\qquad p\in J(s).
\label{eqRbound}
\end{equation}

Write $B(p):=\chi_2'C_2S_1-\chi_1'S_2C_1$ and expand the four factors into exponentials,
$$S_k=\frac i2\left[e^{v_k}e^{-i\theta_k}-e^{-v_k}e^{i\theta_k}\right],\qquad
C_k=\frac 12\left[e^{v_k}e^{-i(\theta_k-\d_k)}+e^{-v_k}e^{i(\theta_k-\d_k)}\right].$$
Then $B$ is a trigonometric polynomial in $p$ with the four frequencies $\mp(t_1+t_2)$, $\pm(t_1-t_2)$:
$$B=\frac i4\left[e^{v_1+v_2}b_1e^{-i(\theta_1+\theta_2)}-e^{v_2-v_1}b_2e^{i(\theta_1-\theta_2)}\right.$$
$$\left.+\,e^{v_1-v_2}b_3e^{-i(\theta_1-\theta_2)}-e^{-v_1-v_2}b_4e^{i(\theta_1+\theta_2)}\right],$$
$$\begin{gathered}b_1=\chi_2'e^{i\d_2}-\chi_1'e^{i\d_1},\qquad b_2=\chi_2'e^{i\d_2}+\chi_1'e^{-i\d_1},\\
b_3=\chi_2'e^{-i\d_2}+\chi_1'e^{i\d_1},\qquad b_4=\chi_2'e^{-i\d_2}-\chi_1'e^{-i\d_1}.\end{gathered}$$
The identities
$$b_2-b_1=2\chi_1'\cos\d_1,\qquad b_3+b_1=2\chi_2'\cos\d_2,$$
together with $|\cos\d_k|=\sqrt\eta+O(\vartheta\eta_0^{-4})$ from \eqref{eqDeltaPin} (and $|\cos\d_2|=\sqrt\eta+O(\epsilon_{2C}\eta_0^{-1/2})$ from Lemma \ref{lemChi} in the free case), give a dichotomy: either $|b_1|\geq\sqrt\eta/2$, or
$$\begin{gathered}|b_2|\ \geq\ 2|\cos\d_1|-|b_1|\ \geq\ \sqrt\eta+O(\vartheta\eta_0^{-4}),\\
|b_3|\ \geq\ 2|\cos\d_2|-|b_1|\ \geq\ \sqrt\eta+O(\vartheta\eta_0^{-4}).\end{gathered}$$
Since the coefficient of $B$ multiplying $b_1$ carries the factor $e^{v_1+v_2}\geq e^{|\nu|}$, and those multiplying $b_2$, $b_3$ carry $e^{\nu}$ and $e^{-\nu}$, in either case the largest of the four coefficients of $B$ is at least $\frac{\sqrt\eta}{8}\,e^{|\nu|}$ in absolute value, once $\d\leq\d_0(C,\eta_0)$ and $C\geq C_1(\eta_0)$ make the $O(\vartheta\eta_0^{-4})$- and $O(\epsilon_{2C}\eta_0^{-1/2})$-terms smaller than $\sqrt\eta/4$; for the latter, $K_{11}e^{-4C+2}\eta_0^{-3/2}\leq\eta_0/4$ once $K_{12}$ is large enough, and the $\d$-part of $\epsilon_c$ is absorbed into $\d_0$.

The pairwise gaps of the four frequencies are $2t_1$, $2t_2$, $2(t_1+t_2)$ and $2(t_2-t_1)$, all at least $t_1=t_2/2$. Since $\left|\int_Ie^{i\omega p}\,dp\right|\leq2/\omega$ for every interval $I$, and $|J(s)|=2C/t$, the mean over $J(s)$ of each cross product of two distinct harmonics is at most $2/C$ times the product of the moduli of their coefficients; hence, with $c_1,\dots,c_4$ denoting the coefficients of $B$,
\begin{multline*}\frac1{|J(s)|}\int_{J(s)}|B|^2\,dp\ \geq\ \sum_j|c_j|^2-\frac 2C\Big(\sum_j|c_j|\Big)^2\ \geq\\ \geq\ \Big(1-\frac{8}C\Big)\max_j|c_j|^2\ \geq\ \frac{\eta}{2^7}\,e^{2|\nu|}\ \geq\ \frac{\eta}{2^7}\,\cosh[2\nu]\end{multline*}
for $C\geq16$; in the free case the terms containing $e^{-v_2}$ are absent and the same dichotomy applies to $b_1$, $b_2$ alone, with the same conclusion in the limiting normalization. For the supremum, each $|b_j|\leq2$, so
\begin{multline*}\sup_J|B|\ \leq\ \frac12\left(e^{v_1+v_2}+e^{\nu}+e^{-\nu}+e^{-v_1-v_2}\right)\ \leq\\ \leq\ 2\,e^{v_1+v_2}\ \leq\ 2\,e^{|\nu|}e^{2\min(v_1,v_2)}\ \leq\ 2\,e^{|\nu|}e^{4C},\end{multline*}
since $\min(v_1,v_2)\leq v_1\leq2C$; hence $\sup_J|B|^2\leq8\cosh[2\nu]\,e^{8C}$.

By \eqref{eqModelB} and \eqref{eqRbound}, since $\g^2(v_1)\g^2(v_2)=4/S$,
$$|b_{t\to2t}(p)|^2=\frac{\left|B(p)+O(\vartheta'e^{4C}\eta_0^{-1/2})\right|^2}{\eta\,S}.$$
For $\d\leq\d_0(C,\eta_0)$ small enough the error term changes the mean and the supremum of $|B|^2$ by relative amounts less than $\tfrac12$ (for the mean, use $\vartheta'e^{4C}\eta_0^{-1/2}\leq\sqrt\eta/2^5$). In the free case the error carries the factor $e^{v_1}$ while the guaranteed coefficient of $B$ may be as small as $\tfrac{\sqrt\eta}4e^{-v_1}$, so the required smallness reads $K\left(\vartheta'+\epsilon_{2C}\right)e^{2v_1}\eta_0^{-1/2}\leq\sqrt\eta/2^5$: since $v_1\leq C$ in the free case, the $\vartheta'$-part follows from $\d\leq\d_0$, and the $\epsilon_c$-part from $C\geq C_1(\eta_0)$ --- $KK_{11}e^{-4C+2}\eta_0^{-3/2}e^{2C}\leq\eta_0/2^5$ once $K_{12}$ is large enough. Thus \eqref{eqSweepMoments} follows.

Finally, the subset of $J(s)$ where $|b_{t\to2t}|^2\geq\frac12\langle|b_{t\to2t}|^2\rangle$ has measure at least
$$\frac{\langle|b|^2\rangle/2}{\sup_J|b|^2}\,|J(s)|\ \geq\ \frac{\cosh[2\nu]/(2^9S)}{2^5\cosh[2\nu]e^{8C}/(\eta S)}\,|J(s)|\ \geq\ 2^{-14}\,\eta_0^2\,e^{-8C}\,|J(s)| ,$$
and on that subset, since
$$\sinh[2v_2]=\sinh[2v_1]\cosh[2\nu]+\cosh[2v_1]\sinh[2\nu]\leq\cosh[2\nu]\,e^{2v_1},$$
$\sinh[2v_1]\leq2v_1e^{2v_1}$ and $v_1\leq2C$,
$$|b_{t\to2t}|^2\ \geq\ \frac{\cosh[2\nu]}{2^9\,S}\ \geq\ \frac1{2^9\sinh[2v_1]e^{2v_1}}\ \geq\ \frac1{2^{10}v_1e^{4v_1}}\ \geq\ \frac{2^{-13}e^{-8C}}{v_1}$$
(in the free case $e^{-2v_1}/(2^{10}\sinh[2v_1])\geq e^{-8C}/(2^{11}v_1)$, with the same conclusion), whence
$\log|a_{t\to2t}|=\frac12\log(1+|b_{t\to2t}|^2)\geq\Phi(v_1,C)$.
\end{proof}

\end{document}
